\label{sec:whole-terminal-control}

Part~II ended with the distinction between strict-subinterval regularity and
the terminal object required by the proof.  The next theorem supplies that
object directly from the datum.  Temporal and spatial depth are fixed
independently.  Because the equation generated every coordinate before either
depth was chosen, all overlapping
rectangles retain the same derivatives of the same history.  The adjacent
rows give terminal traces, and the compatible family reconstructs one common
endpoint from which the equation can restart.

\subsection{The energy-paid base of the producer}

Before the positive rectangles are opened, the kinetic-energy row already
places the first pressure-complete response in two negative Sobolev
topologies.  This is the datum-side base from which the mixed construction
starts.

\begin{proposition}[Energy-paid response faces]
\label{prop:energy-paid-response-faces}
Let \(0<T<T_*\), put \(I_T=(0,T)\), and set
\begin{equation}
 U_0:=\norm{u_0}_2,
 \qquad
 \mathcal E_\nu(T):=
 U_0\left(T+\frac1{2\nu}\right)^{1/2}.
 \label{eq:energy-paid-size}
\end{equation}
With
\[
 \mathcal Q(u,p):=(u\cdot\nabla)u+\nabla p
 =\PP\nabla\cdot(u\otimes u),
\]
the one datum-generated history satisfies
\begin{align}
 \norm u_{L^\infty(I_T;L^2)}
 &\le U_0,
 &
 \norm u_{L^2(I_T;H^1)}
 &\le\mathcal E_\nu(T),
 \label{eq:energy-paid-velocity}\\
 \norm p_{L^{4/3}(I_T;L^2)}
 +\norm{\mathcal Q(u,p)}_{L^{4/3}(I_T;H^{-1})}
 &\le C U_0^{1/2}\mathcal E_\nu(T)^{3/2},
 \label{eq:energy-paid-four-thirds}\\
 \norm p_{L^2(I_T;H^{-1/2})}
 +\norm{\mathcal Q(u,p)}_{L^2(I_T;H^{-3/2})}
 &\le C U_0\mathcal E_\nu(T),
 \label{eq:energy-paid-square}\\
 \norm{u_t}_{L^{4/3}(I_T;H^{-1})}
 &\le
 C U_0^{1/2}\mathcal E_\nu(T)^{3/2}
 {}+\nu T^{1/4}\mathcal E_\nu(T),
 \label{eq:energy-paid-time-four-thirds}\\
 \norm{u_t}_{L^2(I_T;H^{-3/2})}
 &\le
 C U_0\mathcal E_\nu(T)
 {}+\nu\mathcal E_\nu(T).
 \label{eq:energy-paid-time-square}
\end{align}
Every right-hand side is fixed by \(u_0,\nu,T\).
\end{proposition}

\begin{proof}
The energy identity gives
\[
 \sup_{0\le t\le T}\norm{u(t)}_2^2\le U_0^2,
 \qquad
 \int_0^T\norm{\nabla u(t)}_2^2\,dt
 \le\frac{U_0^2}{2\nu}.
\]
Adding the \(L^2\) part of the inhomogeneous \(H^1\) norm proves
\eqref{eq:energy-paid-velocity}.

At each time, the three-dimensional interpolation inequality gives
\[
 \norm{u}_4^2
 \le C\norm u_2^{1/2}\norm{\nabla u}_2^{3/2}.
\]
Since the Riesz transforms are bounded on \(L^2\), while divergence maps
\(L^2\) to \(H^{-1}\), raising this estimate to the power \(4/3\) and
integrating proves \eqref{eq:energy-paid-four-thirds}.

For the second face, duality with
\(H^{1/2}\hookrightarrow L^3\) gives
\[
 \norm{u_i u_j}_{H^{-1/2}}
 \le C\norm u_3^2
 \le C\norm u_2\norm u_{H^1}.
\]
The Riesz formula for pressure and one divergence now give
\eqref{eq:energy-paid-square} after squaring and integrating in time.
Finally,
\[
 u_t=-\nu\PP(-\Delta)u-\mathcal Q(u,p).
\]
The diffusion map sends \(H^1\) to \(H^{-1}\), hence also to
\(H^{-3/2}\).  Its \(L^{4/3}\)-time norm is bounded by
\(T^{1/4}\) times its \(L^2\)-time norm.  Combining these facts with
\eqref{eq:energy-paid-four-thirds} and
\eqref{eq:energy-paid-square} proves the last two estimates.
\end{proof}

These are the first faces of the temporal tower: energy pays the nonlinear
response and \(u_t\) before any positive spatial row is requested.  The
datum-generated theorem below upgrades this same pressure-complete history
to every finite positive adjacent pair on the whole proposed terminal
interval.  The recurrence does not switch histories during that upgrade.

\Needspace{28\baselineskip}
\begin{theorem}[Datum-generated whole-terminal compatible rectangles]
\label{thm:whole-terminal}
Let \(\nu>0\), let \(u_0\in\SSigma(\R^3)\), and let \((u,p)\) be the
pressure-normalized maximal classical solution generated by this datum on
\([0,T_*)\).  If \(T_*<\infty\), set \(I_*=(0,T_*)\) and
\(V_n=\partial_t^n u\).  Then, for every finite \(N,M\in\Nzero\), the
same pressure-complete history produces a whole-terminal rectangle satisfying
\begin{equation}
\begin{aligned}
 \mathcal N_{N,M}(u;I_*)
 :={}&\sum_{n=0}^{N}
 \norm{V_n}_{L^2(I_*;H^{M+1})}^2
 \\
 &+\sum_{n=0}^{N}
 \norm{V_{n+1}}_{L^2(I_*;H^{M-1})}^2
 <\infty.
\end{aligned}
 \label{eq:whole-terminal}
\end{equation}
In adjacent-row coordinates, the output is, for every finite \(n,m\in\Nzero\),
\begin{equation}
 V_n\in L^2(I_*;H^{m+1}),
 \qquad
 V_{n+1}=\partial_tV_n\in L^2(I_*;H^{m-1}).
 \label{eq:datum-adjacent-row}
\end{equation}

Together with the pressure topology in
\eqref{eq:rectangle-pressure-topology}, these bounds place the actual finite
rectangle
\begin{equation*}
 \cR_{N,M}[u_0;T_*]
 :=(V_0,\ldots,V_N;V_{N+1};Q_0,\ldots,Q_N)|_{I_*}
\end{equation*}
in \(\mathscr R_{N,M}(u_0;I_*)\).  If \(N'\le N\) and \(M'\le M\),
restriction satisfies
\begin{equation}
 \rho_{N',M'}^{N,M}\cR_{N,M}[u_0;T_*]
 =\cR_{N',M'}[u_0;T_*].
 \label{eq:restriction}
\end{equation}
Every shared velocity or pressure coordinate in this identity is the same
distributional derivative of the maximal solution.
\end{theorem}

Fix one finite pair \((N,M)\).  Estimate \eqref{eq:whole-terminal} allows the
time cutoff to reach \(T_*\), so every coordinate of that rectangle belongs
to its stated whole-terminal space.  The restriction identity then says that
the corresponding coordinate in any larger rectangle is literally the same
derivative of \((u,p)\).  The estimates arrive on the terminal
interval already attached to one another.

\subsection{Reading the estimate by rows or by rectangles}

The theorem has two useful readings.  The trace argument uses one adjacent
temporal pair at a time; compatibility uses a finite rectangle.  For fixed
\(N\) and \(M\), these are the same finitely many norms.  The passage between
them is worth making explicit because it changes no derivative depth and adds
no infinite-order summation.

Fix \(N,M\) and write
\[
 F_{N,M}(T):=\mathcal N_{N,M}(u;T).
\]
Every term in \(F_{N,M}(T)\) is the integral of a nonnegative function over
\((0,T)\), so \(F_{N,M}\) is increasing.  Monotone convergence gives
\begin{align*}
 \lim_{T\uparrow T_*}F_{N,M}(T)
 &={}
 \sum_{n=0}^{N}
 \norm{V_n}_{L^2(I_*;H^{M+1})}^2
 \\
 &\quad+
 \sum_{n=0}^{N}
 \norm{V_{n+1}}_{L^2(I_*;H^{M-1})}^2.
\end{align*}
There are only finitely many summands.  Applying
\eqref{eq:datum-adjacent-row} to each retained row makes the right side
finite, and restriction to \((0,T)\) gives \eqref{eq:whole-terminal}.

Conversely, fix \(n,m\) and use the rectangle with \((N,M)=(n,m)\).
Estimate \eqref{eq:whole-terminal} bounds the increasing finite-rectangle
norm, so its monotone limit is finite.  The two nonnegative summands
containing \(V_n\) and \(V_{n+1}\) give exactly
\eqref{eq:datum-adjacent-row}.  The temporal and spatial depths remain fixed
throughout both directions.

\section{Pressure completion on the terminal interval}

The upper velocity rows controlled by Theorem~\ref{thm:whole-terminal}
already determine every retained pressure row.  Applying
\eqref{eq:rectangle-pressure-topology} on \(I_*\) gives
\begin{equation}
 \norm{Q_n}_{L^1(I_*;H^M)}
 \le C_{n,M}\sum_{a=0}^{n}
 \norm{V_a}_{L^2(I_*;H^{M+1})}
 \norm{V_{n-a}}_{L^2(I_*;H^{M+1})}
 <\infty.
 \label{eq:whole-terminal-pressure}
\end{equation}
Thus the velocity rectangle itself completes pressure on the terminal
interval.

The same product estimate controls the transport sum
\begin{equation*}
 \sum_{a=0}^{n}\binom na(V_a\cdot\nabla)V_{n-a}
 \quad\hbox{in }L^1(I_*;H^{M-1}).
\end{equation*}
The viscous term belongs to
\(L^2(I_*;H^{M-1})\), and hence to \(L^1(I_*;H^{M-1})\) because
\(T_*<\infty\).  Thus every differentiated momentum row holds in one common
whole-terminal topology.  Because every larger rectangle uses the same
Riesz-transform formula, shared velocity rows determine the same normalized
pressure row under restriction.

\section{Restriction and projective assembly}

The coordinatewise restriction described in Part~II becomes the bonding map
of the whole-terminal finite family.  It deletes higher rows, weakens the
Sobolev topology of the rows that remain, and never changes the retained
functions.

Order the finite pairs by
\[
 (N',M')\preceq(N,M)
 \quad\Longleftrightarrow\quad
 N'\le N\ \hbox{ and }\ M'\le M.
\]
For such a pair, define
\begin{align*}
 \rho_{N',M'}^{N,M}
 &(V_0,\ldots,V_N;V_{N+1};Q_0,\ldots,Q_N)
 \\
 &\qquad:=
 (V_0,\ldots,V_{N'};V_{N'+1};Q_0,\ldots,Q_{N'}),
\end{align*}
where the retained coordinates are carried through the continuous embeddings
\begin{equation*}
 H^{M+1}\hookrightarrow H^{M'+1},
 \qquad
 H^{M-1}\hookrightarrow H^{M'-1},
 \qquad
 H^M\hookrightarrow H^{M'}.
\end{equation*}
This gives a continuous map
\[
 \rho_{N',M'}^{N,M}:
 \mathscr R_{N,M}(u_0;I_*)
 \longrightarrow
 \mathscr R_{N',M'}(u_0;I_*).
\]
When \(N'<N\), the distinguished \(V_{N'+1}\) in the smaller rectangle is an
ordinary retained row of the larger one.  When \(N'=N\), it is the larger
rectangle's distinguished lower-space derivative.  In either case the target
coordinate is the same \(V_{N'+1}\).

The maps obey
\begin{equation*}
 \rho_{N,M}^{N,M}=\operatorname{id},
 \qquad
 \rho_{N'',M''}^{N',M'}\circ\rho_{N',M'}^{N,M}
 =\rho_{N'',M''}^{N,M}
\end{equation*}
whenever
\((N'',M'')\preceq(N',M')\preceq(N,M)\).  Deleting rows in two stages is
the same operation as deleting them at once, and lowering the Sobolev topology
twice leaves the same coordinate as lowering it directly.

The identity and composition laws make the rectangles a projective family.
For the actual solution, compatibility is the concrete identity
\eqref{eq:restriction}: every smaller rectangle is obtained from every larger
one by forgetting coordinates and weakening Sobolev topologies.  No new
velocity or pressure is chosen in that passage.  The differentiated momentum
equations fix the velocity coordinates, and
\eqref{eq:whole-terminal-pressure} fixes the normalized pressure coordinates.

This is the all-orders intersection in its useful form.  It is not one norm of
an infinite array.  It is the assertion that every finite rectangle has its
own whole-terminal bound and that any derivative appearing in two rectangles
is the same function in both.  The first fact supplies depth; the second keeps
identity while the depth increases.

Now fix a temporal row \(V_n\).  Every rectangle tall enough to contain it
uses the same distributional derivative, and increasing \(M\) places that
same function in stronger finite Sobolev spaces.  Their constants may grow
with \(N\) and \(M\), while derivative depth is handled coordinatewise.  The
restriction identities preserve the fact that all of these bounds belong to
one coordinate.

That identity becomes decisive at the terminal time.  An \(H^k\) trace and an
\(H^m\) trace of \(V_n\), with \(k>m\), can define one endpoint only because
both arise from the same preterminal function.  The projective assembly has
already supplied that compatibility before any trace is taken.

\Needspace{14\baselineskip}
\section{What the compatible family gives at all orders}

Continuation begins from the velocity itself.  Fixing the zeroth temporal row
and increasing the spatial depth retains the same \(u\) in every finite
Sobolev space.  For each fixed \(m\), the theorem supplies the velocity
coordinate of a rectangle whose spatial depth is at least \(m\).
Compatibility says that all of these coordinates are the same \(u\).  It
follows that
\begin{equation*}
 u\in\bigcap_{m=0}^{\infty}L^2(I_*;H^m(\R^3)).
\end{equation*}
Every fixed \(L^2_tH^m_x\) norm is finite, with a separate constant permitted
at each spatial order.

The same argument applies to every temporal row.  Fix \(n,q\in\Nzero\) and
choose a rectangle deep enough that its \(n\)-th row is viewed in \(H^q\).
Then
\begin{equation*}
 V_n=\partial_t^n u\in L^2(I_*;H^q)
 \qquad(n,q\in\Nzero).
\end{equation*}
In particular, fix finite \(r,q\in\Nzero\).  The rows
\(V_0,\ldots,V_r\) are the successive distributional time derivatives of
one velocity and all belong to \(L^2(I_*;H^q)\).  Thus
\begin{equation*}
 u\in H^r(I_*;H^q)
 \quad\hbox{for every finite }r,q,
 \qquad
 u\in\bigcap_{r,q\in\Nzero}H^r(I_*;H^q).
\end{equation*}
This is the all-orders velocity realization of the compatible finite family.
Every finite time--space demand is retained separately, and the intersection
is understood coordinatewise.

\subsection{One finite coordinate gives one finite estimate}

The intersection changes the estimate problem because the size of one finite
rectangle can now be carried all the way to any readout contained in it.  Fix
\(N,M\in\Nzero\) and a row \(0\le n\le N\).  Applying the quantitative
adjacent-row trace estimate \eqref{eq:adjacent-trace-quantitative} to \(V_n\)
gives
\begin{align}
 \sup_{0\le t\le T_*}\norm{V_n(t)}_{H^M}^2
 &\le \frac1{T_*}\norm{V_n}_{L^2(I_*;H^M)}^2
 \notag\\
 &\quad+2
 \norm{V_{n+1}}_{L^2(I_*;H^{M-1})}
 \norm{V_n}_{L^2(I_*;H^{M+1})}
 \notag\\
 &\le (1+T_*^{-1})\,
 \mathcal N_{N,M}(u;I_*).
 \label{eq:rectangle-to-trace}
\end{align}
The last line uses \(H^{M+1}\hookrightarrow H^M\) and
\(2xy\le x^2+y^2\).  No all-orders norm enters: the one rectangle on the
right contains both factors that produce the middle trace on the left.

Taking \(N=n\) and \(M=m\) gives, for every fixed \(n,m\in\Nzero\),
\begin{equation}
 V_n\in C([0,T_*];H^m),
 \qquad
 \mathcal M_{n,m}
 :=\sup_{0\le t\le T_*}\norm{V_n(t)}_{H^m}
 \le (1+T_*^{-1})^{1/2}
 \mathcal N_{n,m}(u;I_*)^{1/2}.
 \label{eq:all-row-traces}
\end{equation}
The two-unit gap between the upper and lower rows is doing exact work here:
their shifted pairing lands at the \(H^m\) norm whose endpoint value is being
measured.

Now fix \(k,n\in\Nzero\), \(1\le a\le\infty\), and
\(2\le b\le\infty\).  Choose an integer \(m\) satisfying
\[
 m>k+\frac32-\frac3b.
\]
The spatial embedding \eqref{eq:appendix-Wkp-embedding}, followed by
\eqref{eq:all-row-traces}, gives
\begin{equation}
 \begin{aligned}
  \nabla^k\partial_t^n u
  &\in C([0,T_*];L^b(\R^3))
    \cap L^a(I_*;L^b(\R^3)),\\
  \sup_{0\le t\le T_*}
  \norm{\nabla^k\partial_t^n u(t)}_{L^b}
  &\le C_{m,k,b}(1+T_*^{-1})^{1/2}
  \mathcal N_{n,m}(u;I_*)^{1/2},\\
  \norm{\nabla^k\partial_t^n u}_{L^a(I_*;L^b)}
  &\le C_{m,k,b}\,T_*^{1/a}(1+T_*^{-1})^{1/2}
  \mathcal N_{n,m}(u;I_*)^{1/2},
 \end{aligned}
 \label{eq:finite-projection-mixed-norms}
\end{equation}
where \(T_*^{1/\infty}=1\).  This is the general projection rule.  Once a
finite estimate involving a prescribed number of derivatives and one of these
Sobolev-accessible mixed norms has been named, one sufficiently high finite
coordinate of the compatible family supplies it.
Taking \(a=b=\infty\) after choosing \(m>k+\frac32\) gives the pointwise
version
\[
 \sup_{0\le t\le T_*}
 \norm{\nabla^k\partial_t^n u(t)}_\infty<\infty
 \qquad(n,k\in\Nzero).
\]

Pressure is read from the same coordinates, but now its size can be retained
rather than only its membership.  Fix \(n,s\in\Nzero\) and put
\(r_s:=\max\{s,2\}\).  The normalized identity
\[
 Q_n=R_iR_j\sum_{c=0}^{n}\binom ncV_{c,i}V_{n-c,j}
\]
and the product estimate from Appendix~\ref{app:analytic-framework} give
\begin{align}
 \sup_{0\le t\le T_*}\norm{Q_n(t)}_{H^s}
 &\le C_s\sum_{c=0}^{n}\binom nc
 \mathcal M_{c,r_s}\mathcal M_{n-c,r_s}
 \notag\\
 &\le C_s\,2^n(1+T_*^{-1})
 \mathcal N_{n,r_s}(u;I_*).
 \label{eq:quantitative-pressure-trace}
\end{align}
The common rectangle \((n,r_s)\) contains every factor in the finite
Leibniz sum; the binomial weights contribute exactly \(2^n\).

For fixed \(n,k\in\Nzero\), \(1\le a\le\infty\), and
\(2\le b\le\infty\), choose an integer
\(s>k+\frac32-\frac3b\).  Spatial embedding and
\eqref{eq:quantitative-pressure-trace} then give
\begin{equation}
 \begin{aligned}
 \nabla^kQ_n
 &\in C([0,T_*];L^b(\R^3))
 \cap L^a(I_*;L^b(\R^3)),\\
 \sup_{0\le t\le T_*}\norm{\nabla^kQ_n(t)}_{L^b}
 &\le C_{s,k,b}\,2^n(1+T_*^{-1})
 \mathcal N_{n,r_s}(u;I_*),\\
 \norm{\nabla^kQ_n}_{L^a(I_*;L^b)}
 &\le C_{s,k,b}\,2^nT_*^{1/a}(1+T_*^{-1})
 \mathcal N_{n,r_s}(u;I_*).
 \end{aligned}
 \label{eq:finite-pressure-mixed-norms}
\end{equation}
Equation~\eqref{eq:whole-terminal-pressure} is the direct \(L^1_tH^M_x\)
pressure completion.  Equations~\eqref{eq:quantitative-pressure-trace} and
\eqref{eq:finite-pressure-mixed-norms} are its uniform-time readouts.  Both
come from the same velocity rectangle and the same normalized pressure
formula.

The family is complete before the endpoint is taken.  Its zeroth
row now has a terminal value in every fixed Sobolev space.  Compatibility must
still show that the values obtained from different finite rectangles are one
common endpoint, equal across every continuous Sobolev embedding.  That is the
next step.
